We investigated the deployment challenges for battery SOH prediction that are
absent from small-cell laboratory benchmarks, using a new 18-month dataset of 20
large-format LiFePO$_4$ prismatic cells (50\,Ah) cycled under three operating
conditions.

Its central message is that on large-format cells manufacturing variability, not
algorithm choice, dominates prediction error, and that because one-step accuracy
on their smooth trajectories saturates, models should be judged by the deployment
capabilities they add rather than by one-step rank. We therefore report trivial
baselines and multi-run variance as standard practice.

Guided by this, our hybrid model (MIT LFP pretraining, a convolutional
fingerprint encoder, and MC-Dropout uncertainty) is positioned not as the most
accurate one-step predictor but as the one that supplies what deployment needs:
prediction of a new cell from only its first 30 cycles with no target-cell
training, a substantial accuracy gain once fine tuning begins, and calibrated
uncertainty after a lightweight conformal correction. A ten-scenario
cross-dataset validation shows that fine tuning is the dependable adaptation
mechanism, that zero-shot fingerprint conditioning adds distinctive value when
chemistry and corpus match, and that uncertainty stays calibrated only under
chemistry alignment, a practical guideline for BMS transfer learning. We also
identify pretraining-run variance as a silent reproducibility hazard and
neutralise it with a calibration-cell selection protocol, in a fully
deterministic, byte-for-byte reproducible pipeline.

Building on the same model, we demonstrated two further capabilities aimed
directly at manufacturing variability: an early screen that flags defective cells
from the discharge voltage curve well before the capacity trajectory reveals
them, and a full-coverage probabilistic RUL forecaster that, through a direct
multi-horizon decoder and a 155-cell LFP foundation corpus, forecasts every cell
reaching end of life, including the catastrophic cells recursive baselines
abstain on.

Taken together, this work moves SOH and RUL prediction beyond the small-cell
laboratory setting that has dominated the field, contributing a large-format LFP
benchmark with trivial-baseline controls, evidence that the field's default
evaluation protocol needs rethinking, an uncertainty-calibrated and
deployment-oriented model validated across manufacturers, chemistries, and
operating conditions, and a fully deterministic pipeline with a released
checkpoint. Future work should pursue harder evaluation protocols on which deep
models may separate from trivial baselines, shorter fingerprint windows and
format-matched pretraining, online detection of cells that violate conformal
exchangeability, and pretraining objectives that exclude trivial shortcut
solutions.
